% Preamble template for Tufte-style applied econometrics lecture notes.
% Do not compile this file directly; compile the main notes file instead.
\documentclass{tufte-handout}

\usepackage[T1]{fontenc}
\usepackage[utf8]{inputenc}
\usepackage{ntheorem}
\usepackage{graphicx}
\usepackage{amsmath}
\usepackage{amssymb}
\usepackage{mathtools}
\usepackage{hyperref}
\usepackage{epigraph}
\usepackage{booktabs}
\usepackage{array}
\theoremstyle{break}

\hypersetup{
  colorlinks,
  urlcolor = blue,
  pdfauthor={Swapnil Singh},
  pdfkeywords={econometrics, causal inference, applied econometrics},
  pdftitle={Lecture Notes for Applied Econometrics},
  pdfpagemode=UseNone
}

\newtheorem{ruleN}{Rule}
\newtheorem{thmN}{Theorem}
\newtheorem{assN}{Assumption}
\newtheorem{defN}{Definition}
\newtheorem{exmp}{Example}
\newtheorem{cmt}{Comment}
\newtheorem{proof}{Proof}

\newcommand{\continuation}{??}
\newtheorem*{excont}{Example \continuation}
\newenvironment{continueexample}[1]
 {\renewcommand{\continuation}{\ref{#1}}\excont[continued]}
 {\endexcont}

\newcommand{\bY}{\mathbf{Y}}
\newcommand{\bX}{\mathbf{X}}
\newcommand{\bD}{\mathbf{D}}
\newcommand{\E}{\mathbb{E}}
\newcommand{\Prob}{\mathbb{P}}
\newcommand{\Var}{\mathrm{Var}}
\newcommand{\Cov}{\mathrm{Cov}}
\newcommand{\se}{\mathrm{SE}}
\newcommand{\CI}{\mathrm{CI}}
\newcommand{\ATE}{\mathrm{ATE}}
\newcommand{\ATT}{\mathrm{ATT}}
\newcommand{\TOT}{\mathrm{TOT}}
\newcommand{\LATE}{\mathrm{LATE}}
\newcommand{\ITT}{\mathrm{ITT}}
\newcommand{\RD}{\mathrm{RD}}
\newcommand{\DiD}{\mathrm{DiD}}
\newcommand{\MMSE}{\mathrm{MMSE}}
\newcommand{\OLS}{\mathrm{OLS}}
\newcommand{\Normal}{\mathcal{N}}
\newcommand{\plim}{\operatorname*{plim}}
\DeclareMathOperator*{\argmin}{arg\,min}
\DeclareMathOperator{\Supp}{Supp}

\newcommand\independent{\protect\mathpalette{\protect\independenT}{\perp}}
\def\independenT#1#2{\mathrel{\rlap{$#1#2$}\mkern2mu{#1#2}}}
\newcommand{\indep}{\perp\!\!\!\perp}

\usepackage{cleveref}
\crefname{appsec}{appendix}{appendices}
\crefname{appsubsec}{appendix}{appendices}
\crefname{assumption}{assumption}{assumptions}
\crefname{equation}{equation}{equations}
\crefname{exmp}{example}{examples}
\crefname{assN}{assumption}{assumptions}
\crefname{cmt}{comment}{comments}
\crefname{defN}{definition}{definitions}

\usepackage[nolist]{acronym}
\begin{acronym}
  \acro{ATE}{average treatment effect}
  \acro{ATT}{average treatment effect on the treated}
  \acro{CIA}{conditional independence assumption}
  \acro{CI}{confidence interval}
  \acro{CLT}{central limit theorem}
  \acro{DiD}{difference-in-differences}
  \acro{FE}{fixed effects}
  \acro{FWL}{Frisch-Waugh-Lovell}
  \acro{IV}{instrumental variables}
  \acro{LATE}{local average treatment effect}
  \acro{OLS}{ordinary least squares}
  \acro{OVB}{omitted-variable bias}
  \acro{PO}{potential outcomes}
  \acro{RCT}{randomized controlled trial}
  \acro{RD}{regression discontinuity}
  \acro{SUTVA}{stable unit treatment value assignment}
  \acro{TWFE}{two-way fixed effects}
\end{acronym}

\def\inprobHIGH{\,{\buildrel p \over \rightarrow}\,}
\def\inprob{\,{\inprobHIGH}\,}
\def\indistHIGH{\,{\buildrel d \over \rightarrow}\,}
\def\indist{\,{\indistHIGH}\,}

\usepackage[many]{tcolorbox}
\definecolor{main}{HTML}{3F6EA7}
\definecolor{sub}{HTML}{F2F7FD}
\definecolor{sub2}{HTML}{FFF5E8}
\definecolor{mainline}{HTML}{D3DFEE}
\definecolor{warn}{HTML}{C7771B}
\definecolor{warnline}{HTML}{EAD9C0}

\tcbset{
    enhanced,
    breakable,
    colback = white,
    boxrule = 0.5pt,
    arc = 2pt,
    left = 9pt,
    right = 9pt,
    top = 7pt,
    bottom = 7pt,
    before skip = 0.3cm,
    after skip = 0.45cm
}

\newtcolorbox{boxD}{
    colback = sub,
    colframe = mainline,
    borderline west = {2.2pt}{0pt}{main},
    borderline north = {0.6pt}{0pt}{main!25},
    borderline south = {0.6pt}{0pt}{main!15}
}

\newtcolorbox{boxF}{
    colback = sub2,
    colframe = warnline,
    borderline west = {2.2pt}{0pt}{warn},
    borderline north = {0.6pt}{0pt}{warn!25},
    borderline south = {0.6pt}{0pt}{warn!15}
}

\usepackage{colortbl}
\usepackage{tikz}
\usepackage{verbatim}
\usetikzlibrary{positioning}
\usetikzlibrary{snakes}
\usetikzlibrary{calc}
\usetikzlibrary{arrows}
\usetikzlibrary{decorations.markings}
\usetikzlibrary{shapes.misc}
\usetikzlibrary{matrix,shapes,arrows,fit,tikzmark}


\title{Applied Econometrics: Session 3}
\author{PhD Intensive Course}
\date{}

\hypersetup{
  pdftitle={Applied Econometrics Session 3: Stats Review},
  pdfauthor={Swapnil Singh},
  pdfkeywords={econometrics, inference, standard errors, confidence intervals, p-values, causal inference}
}

\setlength{\headheight}{26pt}

\begin{document}
\maketitle

\begin{abstract}
This note reviews the statistical layer that sits on top of causal design.
The main theme is the ``most dangerous equation'': the standard error of an
average falls at rate $1/\sqrt{n}$, not at rate $1/n$. This simple fact explains
why small groups often look extreme in rankings, why randomized experiments
can be correct in design but noisy in finite samples, and why applied
econometrics must report uncertainty together with point estimates. We cover
estimands, estimators, estimates, sampling variation, standard errors,
confidence intervals, hypothesis tests, p-values, power, multiple testing, and
the relationship between uncertainty and causal interpretation.
\end{abstract}

\begin{boxD}
\textbf{Reading and objective.} Read Chapter 03, ``Stats Review: The Most
Dangerous Equation,'' from \emph{Causal Inference for the Brave and True}.
The objective is not to memorize formulas mechanically. The objective is to
learn how to read an empirical result as one noisy realization from a sampling
process, and to separate three questions: what is the causal target, what is
the estimator, and how uncertain is the estimate?
\end{boxD}

\section{Why This Lecture Matters}

The first two lectures were about causal meaning. We asked what it means to
say that a treatment caused an outcome, how potential outcomes define causal
effects, why selection bias is dangerous, and why regression is best understood
as a conditional expectation approximation unless stronger assumptions give it
causal content.

This lecture adds the statistical layer. Even when the causal design is
excellent, the estimate in a paper is still computed from a finite sample. A
different random sample, or a different random assignment, would usually have
produced a different number. That movement from sample to sample is called
sampling variation.

\begin{boxF}
\textbf{Key warning.} A credible research design can produce an imprecise
estimate. A precise estimate can still be biased. Identification and inference
solve different problems.
\end{boxF}

An applied empirical claim should therefore contain at least four parts:
\begin{enumerate}
  \item the target parameter or estimand;
  \item the estimator used to learn about that target;
  \item the numerical estimate obtained in the sample;
  \item a measure of uncertainty, usually a standard error or confidence
  interval.
\end{enumerate}

\section{Estimand, Estimator, and Estimate}

\begin{defN}[Estimand]
An estimand is the population object we want to learn. In causal work it is
often an average treatment effect,
\[
  \ATE = \E[Y_1 - Y_0],
\]
or an average treatment effect for a subpopulation, such as $\ATT$.
\end{defN}

\begin{defN}[Estimator]
An estimator is a rule that maps data into a number. For example, in a
randomized experiment with a treated group and a control group, a natural
estimator of the average treatment effect is
\[
  \hat{\tau} = \bar{Y}_1 - \bar{Y}_0.
\]
Before the data are realized, $\hat{\tau}$ is a random variable.
\end{defN}

\begin{defN}[Estimate]
An estimate is the actual numerical value produced by applying an estimator to
one dataset. If the treated mean is $76.3$ and the control mean is $72.1$, the
estimate is
\[
  \hat{\tau} = 76.3 - 72.1 = 4.2.
\]
\end{defN}

Students often mix these objects. The estimand is the thing in the world we
would like to know. The estimator is the recipe. The estimate is the number on
the page. Standard errors describe the behavior of the recipe across repeated
samples, not the moral importance of the number and not the truth of the
identifying assumption.

\begin{exmp}[Tutoring program]
Suppose a university randomly offers tutoring to half of first-year students.
The causal estimand might be the average effect of tutoring on exam scores for
this population. The estimator might be the treated mean minus the control
mean. The estimate might be $4.2$ exam points. The standard error then tells us
how much this difference-in-means estimator would typically move if the
experiment were repeated.
\end{exmp}

\section{Sampling Variation}

Imagine a large population of students with average exam score $\mu=70$ and
individual standard deviation $\sigma=12$. If we draw one classroom of $25$
students, the classroom average might be $66$. If we draw another classroom of
$25$ students, the average might be $73$. Neither classroom average is the
population mean. Each is one noisy estimate of the population mean.

\begin{defN}[Sampling variation]
Sampling variation is the variation in an estimator across repeated samples
generated by the same sampling process.
\end{defN}

\begin{defN}[Sampling distribution]
The sampling distribution of an estimator is the probability distribution of
the estimator across repeated samples. The population distribution describes
individual observations. The sampling distribution describes statistics
computed from samples.
\end{defN}

This distinction matters. A standard deviation describes the dispersion of
individual outcomes. A standard error describes the dispersion of an estimator.
For instance, students may differ a lot in exam scores, but if we observe a
very large number of students, the average exam score may still be estimated
very precisely.

\section{The School-Ranking Trap}

The Python Causality Handbook chapter begins with a simple but powerful idea:
small groups are more likely to produce extreme averages. This is the logic
behind Howard Wainer's phrase ``the most dangerous equation.''

Consider school rankings based on average test scores. Some schools have only
20 tested students, while others have 500. Even if all schools were equally
good, the small schools would have much noisier average scores because the
average is based on fewer observations. Therefore small schools would be
overrepresented among the best observed schools. They would also be
overrepresented among the worst observed schools.

\begin{boxD}
\textbf{Diagnostic for rankings.} If a small group appears surprisingly often
at the top of a ranking, check the bottom of the ranking. If the same kind of
small group also appears at the bottom, the pattern is probably about noise,
not excellence.
\end{boxD}

\begin{exmp}[Schools with no size effect]
Suppose school quality is independent of enrollment. Let the true school mean
be
\[
  \mu_j = 70 + q_j,
  \qquad q_j \sim \Normal(0,2^2),
\]
and let the observed school average be
\[
  \bar{Y}_j = \mu_j + u_j,
  \qquad u_j \sim \Normal\left(0,\frac{12^2}{n_j}\right).
\]
The variance of the noise term is larger when $n_j$ is smaller. Thus the
smallest schools naturally generate the most volatile observed averages.
\end{exmp}

The course script \texttt{lectures/code/lecture-3-stats-review.py} simulates
this exact logic. It creates many schools with different enrollments, assigns
true school quality independently of size, and then adds sampling noise with
standard deviation $12/\sqrt{n_j}$. The script prints the average school size
among top, middle, and bottom observed performers. The intended lesson is that
extreme observed performance need not reveal extreme true quality.

\section{The Most Dangerous Equation}

\begin{thmN}[Standard error of the sample mean]
Let $Y_1,\ldots,Y_n$ be independent observations with
\[
  \E[Y_i]=\mu,
  \qquad
  \Var(Y_i)=\sigma^2.
\]
The sample mean
\[
  \bar{Y} = \frac{1}{n}\sum_{i=1}^n Y_i
\]
has variance
\[
  \Var(\bar{Y}) = \frac{\sigma^2}{n}
\]
and standard error
\[
  \se(\bar{Y}) = \frac{\sigma}{\sqrt{n}}.
\]
\end{thmN}

\begin{proof}
Using independence,
\[
  \Var(\bar{Y})
  =
  \Var\left(\frac{1}{n}\sum_{i=1}^n Y_i\right)
  =
  \frac{1}{n^2}\Var\left(\sum_{i=1}^n Y_i\right)
  =
  \frac{1}{n^2}\sum_{i=1}^n \Var(Y_i).
\]
Since each variance is $\sigma^2$,
\[
  \Var(\bar{Y})
  =
  \frac{1}{n^2}n\sigma^2
  =
  \frac{\sigma^2}{n}.
\]
Taking the square root gives $\se(\bar{Y})=\sigma/\sqrt{n}$.
\end{proof}

The dangerous part is the square root. Precision improves slowly. To cut the
standard error in half, we need four times as many observations.

\begin{center}
\begin{tabular}{rcc}
\toprule
Sample size $n$ & $\se(\bar{Y})=12/\sqrt{n}$ & Interpretation \\
\midrule
25 & 2.40 & noisy average \\
100 & 1.20 & half as noisy \\
400 & 0.60 & half again \\
1600 & 0.30 & half again \\
\bottomrule
\end{tabular}
\end{center}

\begin{boxF}
\textbf{Common mistake.} Do not say that increasing the sample from $100$ to
$200$ cuts the standard error in half. It cuts the standard error by a factor
of $\sqrt{100/200}=1/\sqrt{2}$, about $0.71$.
\end{boxF}

\section{Standard Deviation Versus Standard Error}

\begin{center}
\begin{tabular}{p{0.42\linewidth}p{0.42\linewidth}}
\toprule
Standard deviation & Standard error \\
\midrule
Spread of individual outcomes & Spread of an estimator \\
Describes students, workers, firms, schools & Describes $\bar{Y}$, $\hat{\beta}$, $\widehat{\ATE}$ \\
Usually does not vanish with $n$ & Usually shrinks as sample size grows \\
Answers: how different are units? & Answers: how precise is the estimate? \\
\bottomrule
\end{tabular}
\end{center}

If students differ greatly in ability, the standard deviation of scores may be
large. But if we observe thousands of students, the average score may be known
quite precisely. Conversely, a small pilot study may have an average that moves
a lot across samples even when individual outcomes are not especially
variable.

\section{Estimators Are Random Variables}

Before data are observed, the sample mean is random:
\[
  \bar{Y} = \frac{1}{n}\sum_{i=1}^n Y_i.
\]
After data are observed, $\bar{Y}$ becomes a number. In inference, we reason
about the estimator before the data were realized. This is why we can say that
a reported estimate is one point from a cloud of estimates that could have
been obtained under repeated sampling.

\begin{assN}[Independent sampling]
For the elementary formulas in this lecture, observations are treated as
independent draws from a common population distribution. Later in the course,
we will relax this with heteroskedasticity-robust, clustered, and time-series
standard errors.
\end{assN}

\begin{boxD}
\textbf{Mental model.} The estimate in a table is not the whole cloud. It is
one draw from the cloud. The standard error estimates the cloud's width.
\end{boxD}

\section{Bias and Variance}

It is useful to separate two components of estimation error:
\[
  \hat{\theta}-\theta
  =
  \underbrace{\E[\hat{\theta}]-\theta}_{\text{bias}}
  +
  \underbrace{\hat{\theta}-\E[\hat{\theta}]}_{\text{sampling noise}}.
\]
Identification assumptions aim at the first term. Standard errors describe
the second term.

\begin{boxF}
\textbf{Bad inference habit.} A small standard error does not rescue a biased
research design. It only says that the biased estimator is tightly estimated.
\end{boxF}

For causal work this distinction is essential. A randomized experiment can
make $\E[\hat{\tau}]$ equal to the causal estimand, but in a small experiment
$\hat{\tau}$ can still be far from the truth. An observational regression with
omitted variables can have a tiny standard error and still estimate the wrong
object.

\section{Standard Error of a Mean}

In practice, the population standard deviation $\sigma$ is unknown. We estimate
it with the sample standard deviation
\[
  s
  =
  \sqrt{
    \frac{1}{n-1}\sum_{i=1}^n (Y_i-\bar{Y})^2
  }.
\]
Then the estimated standard error of the mean is
\[
  \widehat{\se}(\bar{Y})=\frac{s}{\sqrt{n}}.
\]

The denominator $n-1$ appears because $\bar{Y}$ is estimated from the same data.
Once the deviations $Y_i-\bar{Y}$ are formed, they must add to zero, so only
$n-1$ deviations are free to vary independently.

\begin{exmp}[Exam scores]
Suppose a class has $\bar{Y}=72$, $s=12$, and $n=25$. Then
\[
  \widehat{\se}(\bar{Y}) = \frac{12}{\sqrt{25}} = 2.40.
\]
A rough 95 percent margin of error is
\[
  1.96 \times 2.40 = 4.70.
\]
With $n=100$ and the same standard deviation, the standard error is $1.20$ and
the rough margin is $2.35$.
\end{exmp}

\section{Treatment-Control Differences}

In a randomized experiment with binary treatment $D_i$, define
\[
  \bar{Y}_1 = \frac{1}{n_1}\sum_{D_i=1}Y_i,
  \qquad
  \bar{Y}_0 = \frac{1}{n_0}\sum_{D_i=0}Y_i.
\]
The difference-in-means estimator is
\[
  \widehat{\ATE} = \bar{Y}_1-\bar{Y}_0.
\]
Random assignment gives the causal interpretation; the standard error gives
the precision.

\begin{thmN}[Standard error of a difference in means]
If treatment and control sample means are independent, then
\[
  \Var(\bar{Y}_1-\bar{Y}_0)
  =
  \frac{\sigma_1^2}{n_1}
  +
  \frac{\sigma_0^2}{n_0}.
\]
The usual plug-in standard error is
\[
  \widehat{\se}(\bar{Y}_1-\bar{Y}_0)
  =
  \sqrt{
    \frac{s_1^2}{n_1}
    +
    \frac{s_0^2}{n_0}
  }.
\]
\end{thmN}

\begin{proof}
Use
\[
  \Var(A-B)=\Var(A)+\Var(B)-2\Cov(A,B).
\]
When the two sample means are independent, the covariance term is zero. The
variance of each group mean is the group variance divided by the group sample
size. Substitution gives the result.
\end{proof}

\begin{exmp}[Balanced and unbalanced experiments]
Suppose $s_1=s_0=12$ and the total sample size is $200$. If $n_1=n_0=100$,
\[
  \widehat{\se}
  =
  \sqrt{\frac{12^2}{100}+\frac{12^2}{100}}
  =
  1.70.
\]
If instead $n_1=20$ and $n_0=180$,
\[
  \widehat{\se}
  =
  \sqrt{\frac{12^2}{20}+\frac{12^2}{180}}
  =
  2.83.
\]
For equal variances and a fixed total sample size, balanced groups are most
precise.
\end{exmp}

\section{The Central Limit Theorem}

The central limit theorem is the reason that normal approximations appear so
often in applied econometrics.

\begin{thmN}[Central limit theorem, informal]
If $Y_1,\ldots,Y_n$ are independent draws with mean $\mu$ and finite variance
$\sigma^2$, then for large $n$,
\[
  \frac{\bar{Y}-\mu}{\sigma/\sqrt{n}}
  \approx
  \Normal(0,1).
\]
\end{thmN}

The CLT does not say that individual outcomes are normal. It says that averages
are approximately normal in large samples. This is why we can use normal
critical values for sample means and many regression estimators even when the
raw data are skewed.

\begin{boxF}
\textbf{Regularity matters.} The CLT can work poorly with very small samples,
extreme outliers, heavy tails, strong dependence, or badly behaved sampling
schemes. The formula is a workhorse, not a license to ignore the data
generating process.
\end{boxF}

\section{Confidence Intervals}

For an estimate $\hat{\theta}$ with estimated standard error
$\widehat{\se}(\hat{\theta})$, the usual large-sample 95 percent confidence
interval is
\[
  \hat{\theta}
  \pm
  1.96\,\widehat{\se}(\hat{\theta}).
\]
The center is the estimate. The radius is about two standard errors. Wider
intervals mean less precision.

\begin{exmp}[Tutoring effect]
Suppose a randomized tutoring program gives
\[
  \hat{\tau}=4.0,
  \qquad
  \widehat{\se}(\hat{\tau})=1.5.
\]
The 95 percent confidence interval is
\[
  4.0 \pm 1.96(1.5)
  =
  [1.06,6.94].
\]
The data are compatible with a modest positive effect and a large positive
effect. Under the usual 95 percent rule, zero is not in the interval.
\end{exmp}

\subsection{Where the Formula Comes From}

The CLT gives
\[
  Z=
  \frac{\hat{\theta}-\theta}{\se(\hat{\theta})}
  \approx \Normal(0,1).
\]
For a standard normal variable,
\[
  \Prob(-1.96 \leq Z \leq 1.96) \approx 0.95.
\]
Substitute the statistic:
\[
  \Prob\left(
    -1.96
    \leq
    \frac{\hat{\theta}-\theta}{\se(\hat{\theta})}
    \leq
    1.96
  \right)
  \approx 0.95.
\]
Rearranging the inequalities gives
\[
  \Prob\left(
    \hat{\theta}-1.96\se(\hat{\theta})
    \leq
    \theta
    \leq
    \hat{\theta}+1.96\se(\hat{\theta})
  \right)
  \approx 0.95.
\]
In practice we replace the unknown standard error with an estimated standard
error.

\subsection{What 95 Percent Confidence Means}

Frequentist confidence is about the procedure. If we repeated the study many
times and computed an interval each time using the same formula, about 95
percent of those intervals would cover the true parameter.

After seeing one interval, the parameter is fixed. The interval either covers
it or it does not. Strictly speaking, it is not correct to say that there is a
95 percent probability that this exact interval contains the true parameter.

\begin{boxD}
\textbf{Good language.} ``Using a procedure that has about 95 percent
coverage in repeated samples, we obtain the interval $[1.06,6.94]$.''
\end{boxD}

\subsection{Do Not Compare Two Separate Intervals}

To test whether treatment and control means differ, use the standard error for
the difference:
\[
  (\bar{Y}_1-\bar{Y}_0)
  \pm
  1.96
  \sqrt{
    \frac{s_1^2}{n_1}
    +
    \frac{s_0^2}{n_0}
  }.
\]
Separate confidence intervals for $\bar{Y}_1$ and $\bar{Y}_0$ answer separate
questions. Their visual overlap is not the correct test of the treatment
effect. The uncertainty must be attached to the estimator being interpreted.

\section{Hypothesis Tests}

A hypothesis test starts by assuming a benchmark value is true:
\[
  H_0:\theta=\theta_0.
\]
Then it asks whether the observed estimate is surprising under that benchmark.
For a large-sample normal test, the test statistic is
\[
  z=
  \frac{\hat{\theta}-\theta_0}
       {\widehat{\se}(\hat{\theta})}.
\]
The numerator is distance from the null. The denominator is typical sampling
noise. A value of $z=2$ means that the estimate is about two standard errors
away from the null.

\begin{boxD}
\textbf{Two-sided 5 percent rule.} Reject $H_0$ at the 5 percent level if
$|z|>1.96$.
\end{boxD}

\begin{exmp}[Testing a tutoring effect]
Let $\hat{\tau}=4.0$, $\widehat{\se}=1.5$, and test $H_0:\tau=0$. Then
\[
  z=\frac{4.0-0}{1.5}=2.67.
\]
Since $2.67>1.96$, we reject zero at the 5 percent level in a two-sided test.
\end{exmp}

Do not confuse rejecting a null with proving a theory. The test is conditional
on assumptions about sampling, standard errors, and the null. It is one part
of empirical reasoning, not the whole argument.

\section{P-values}

\begin{defN}[P-value]
The p-value is the probability, under the null hypothesis, of observing a test
statistic at least as extreme as the one actually observed.
\end{defN}

For a two-sided normal test,
\[
  p = 2[1-\Phi(|z|)],
\]
where $\Phi$ is the standard normal cumulative distribution function.

The most common mistake is to read a p-value as
\[
  \Prob(H_0 \mid \text{data}).
\]
That is not what it is. A p-value is closer to
\[
  \Prob(\text{data as or more extreme} \mid H_0).
\]

\begin{boxF}
\textbf{Plain English.} A small p-value says the estimate would be unusual if
the null were true. It does not say the probability that the null is true.
\end{boxF}

\subsection{One-Sided and Two-Sided Tests}

A two-sided alternative is
\[
  H_1:\theta\neq\theta_0.
\]
It rejects for estimates far above or far below the null. A one-sided
alternative is
\[
  H_1:\theta>\theta_0
  \quad\text{or}\quad
  H_1:\theta<\theta_0.
\]
One-sided tests should be chosen before looking at the data and should be
justified by the research question. In applied work, two-sided tests are the
default.

\section{Statistical Significance and Economic Importance}

Statistical significance is not economic importance. With a huge dataset, a
tiny effect can be statistically significant. With a small dataset, a large
effect can be statistically imprecise.

\begin{center}
\begin{tabular}{lcc}
\toprule
Result & Estimate & Standard error \\
\midrule
Huge dataset & 0.03 & 0.01 \\
Small study & 4.00 & 2.50 \\
\bottomrule
\end{tabular}
\end{center}

The first result gives $z=3$, so it is statistically significant at common
levels, but an effect of $0.03$ may be economically irrelevant. The second
result gives $z=1.6$, so it is not significant at the 5 percent level, but an
effect of $4$ units may matter a lot. Applied interpretation must discuss
magnitude, uncertainty, and design.

\section{Failure to Reject Is Not Evidence of Zero}

Suppose
\[
  \hat{\theta}=4.0,
  \qquad
  \widehat{\se}=3.0.
\]
Testing $H_0:\theta=0$ gives
\[
  z=\frac{4.0}{3.0}=1.33.
\]
The 95 percent confidence interval is
\[
  4.0 \pm 1.96(3.0) = [-1.88,9.88].
\]
We cannot reject zero at the 5 percent level. But the interval also contains
large positive effects. A better interpretation is: the estimate is too
imprecise to distinguish zero from meaningful positive effects.

\begin{boxF}
\textbf{Bad sentence.} ``The treatment had no effect because the result is not
statistically significant.''

\textbf{Better sentence.} ``The study is too imprecise to rule out either zero
or economically meaningful positive effects.''
\end{boxF}

\section{Power}

\begin{defN}[Power]
Power is the probability of rejecting the null hypothesis when a specific
alternative is true.
\end{defN}

Power depends on the true effect size, the standard error, the significance
level, and the form of the test. For a two-sided test of $H_0:\theta=0$, the
test rejects when
\[
  |\hat{\theta}/\se(\hat{\theta})|>1.96.
\]
If the true effect is small relative to the standard error, the test will
reject rarely. If the true effect is large relative to the standard error, the
test will reject often.

\begin{exmp}[Approximate power logic]
Suppose the true effect is $\tau=4$ and the standard error is approximately
$2.4$. Then the estimator is roughly
\[
  \hat{\tau}\sim \Normal(4,2.4^2).
\]
For a two-sided 5 percent test of zero, rejection requires
$|\hat{\tau}|>1.96(2.4)=4.70$. Since the true mean of the sampling distribution
is only $4$, many samples will fail to reject even though the effect is truly
positive. The study is underpowered.
\end{exmp}

Low-power studies have two bad features. First, they often miss real effects.
Second, when they do reject, the significant estimates are often exaggerated in
magnitude because only unusually large draws cross the significance threshold.

\section{Multiple Testing}

If one test is run at the 5 percent level and the null is true, the probability
of a false rejection is 5 percent. If many tests are run, the chance of at
least one false rejection can become large.

If $m$ independent true null hypotheses are tested at level $\alpha$, the
probability of at least one false positive is
\[
  1-(1-\alpha)^m.
\]
With $\alpha=0.05$ and $m=20$,
\[
  1-0.95^{20}\approx 0.64.
\]
Thus, if we search across many outcomes, subgroups, and specifications, some
small p-values can appear by chance.

\begin{boxD}
\textbf{Practical discipline.} Pre-specify main outcomes when possible. Report
the family of results, not only the most impressive one. Label exploratory
analysis as exploratory. Consider corrections such as Bonferroni, Holm, or
false-discovery-rate procedures when the research design involves many tests.
\end{boxD}

\section{Uncertainty and Causal Estimates}

Causal interpretation and statistical precision answer different questions.

\begin{center}
\begin{tabular}{p{0.36\linewidth}p{0.50\linewidth}}
\toprule
Problem & Main question \\
\midrule
Identification & Is the estimator centered on the causal estimand? \\
Inference & How dispersed is the estimator around that center? \\
Measurement & Are treatment and outcome recorded correctly? \\
External validity & Does this population answer the policy question? \\
\bottomrule
\end{tabular}
\end{center}

In a randomized experiment,
\[
  D_i \independent (Y_{1i},Y_{0i})
\]
solves the selection problem in expectation. But the finite-sample difference
in means still varies across random samples and random assignments:
\[
  \hat{\tau}=\bar{Y}_1-\bar{Y}_0.
\]
Randomization removes systematic selection bias; it does not remove luck.

In observational designs, the standard error formula must match the sampling
process and dependence structure. Heteroskedasticity, clustering, serial
correlation, generated regressors, attrition, and noncompliance can all change
the appropriate inference method.

\begin{boxF}
\textbf{Course habit.} Never use the word ``causal'' because a p-value is
small. Use the word ``causal'' because the identifying variation and
assumptions justify a causal interpretation.
\end{boxF}

\section{Using the Python Example}

The script
\[
  \texttt{lectures/code/lecture-3-stats-review.py}
\]
is a companion to this lecture. From the repository root, run
\[
  \texttt{python lectures/code/lecture-3-stats-review.py}
\]
or, to save figures,
\[
  \texttt{python lectures/code/lecture-3-stats-review.py --save-plots}.
\]

The script contains two demonstrations.
\begin{enumerate}
  \item \textbf{School rankings.} School size has no causal effect in the data
  generating process, but small schools appear in both tails of observed
  rankings because their averages have larger standard errors.
  \item \textbf{Repeated experiments.} The true treatment effect is fixed, but
  repeated randomized studies produce different estimates. As sample size
  grows, the sampling distribution tightens around the truth.
\end{enumerate}

When reading the output, focus on three quantities: the point estimate, the
standard error, and the confidence interval. Then ask whether the interval is
narrow enough to answer the substantive question.

\section{How to Read a Results Table}

When reading any empirical results table, do not start with stars. Read in this
order:
\begin{enumerate}
  \item What is the estimand: ATE, ATT, regression slope, reduced form, or
  something else?
  \item What identifying assumption gives the estimate causal meaning?
  \item What are the units and magnitude of the point estimate?
  \item What is the standard error or confidence interval?
  \item Is the estimate precise enough to answer the economic question?
  \item Are there multiple outcomes, subgroups, or specifications that change
  how we should read the p-values?
\end{enumerate}

A useful writing template is:
\begin{quote}
We estimate that $X$ changes $Y$ by $\hat{\theta}$ units
$(\widehat{\se}=s)$, with a 95 percent confidence interval of $[L,U]$. The
identifying variation comes from \ldots. The estimate is economically
large/small because \ldots. The main threat to causal interpretation is
\ldots.
\end{quote}

\section{Worked Problems}

\subsection{Problem 1: Standard Error of a Mean}

A researcher observes $n=64$ workers. Their average weekly hours are
$\bar{Y}=38$ and the sample standard deviation is $s=8$.

\textbf{Solution.} The estimated standard error is
\[
  \widehat{\se}(\bar{Y})=\frac{8}{\sqrt{64}}=1.
\]
The approximate 95 percent confidence interval is
\[
  38 \pm 1.96(1) = [36.04,39.96].
\]
Interpretation: the average is estimated with a margin of about two hours.

\subsection{Problem 2: Difference in Means}

A randomized training program has $n_1=100$ treated workers and $n_0=100$
control workers. Mean monthly earnings are $\bar{Y}_1=2200$ and
$\bar{Y}_0=2100$. The sample standard deviations are $s_1=500$ and $s_0=600$.

\textbf{Solution.} The estimated effect is
\[
  \hat{\tau}=2200-2100=100.
\]
The standard error is
\[
  \widehat{\se}(\hat{\tau})
  =
  \sqrt{\frac{500^2}{100}+\frac{600^2}{100}}
  =
  \sqrt{2500+3600}
  =
  78.10.
\]
The 95 percent confidence interval is
\[
  100 \pm 1.96(78.10)
  =
  [-53.08,253.08].
\]
This experiment is compatible with a negative effect, zero, and a moderately
large positive effect. The point estimate is positive, but precision is weak.

\subsection{Problem 3: Hypothesis Test and P-value}

Using Problem 2, test $H_0:\tau=0$ against a two-sided alternative.

\textbf{Solution.} The test statistic is
\[
  z=\frac{100}{78.10}=1.28.
\]
Since $1.28<1.96$, we do not reject zero at the 5 percent level. The two-sided
p-value is
\[
  p=2[1-\Phi(1.28)]\approx 0.20.
\]
This does not prove the effect is zero. It says that an estimate this far from
zero is not very unusual if the null were true.

\subsection{Problem 4: Multiple Testing}

A paper tests 10 independent outcomes at the 5 percent level, and suppose all
10 null hypotheses are true. What is the probability of at least one false
positive?

\textbf{Solution.}
\[
  1-(1-0.05)^{10}
  =
  1-0.95^{10}
  \approx
  0.40.
\]
Even with no true effects, there is about a 40 percent chance of at least one
significant result. This is why selective reporting is dangerous.

\subsection{Problem 5: Causal Meaning Versus Precision}

An observational regression estimates that private tutoring raises scores by
5 points with standard error 0.5. Should we conclude that tutoring causes
scores to rise by 5 points?

\textbf{Solution.} Not from precision alone. The estimate is precise: a 95
percent interval is roughly $[4.02,5.98]$. But causal interpretation requires
an identifying assumption. Students who receive tutoring may differ in family
income, motivation, prior preparation, or school quality. The small standard
error only says the regression coefficient is tightly estimated; it does not
solve selection bias.

\section{Checklist for Students}

\begin{enumerate}
  \item Can I state the estimand in words and in notation?
  \item Can I distinguish the estimator from the estimate?
  \item Do I know what sampling process makes the estimator random?
  \item Did I report the standard error in the same units as the estimate?
  \item Did I compute the confidence interval for the estimator I interpret?
  \item Did I avoid comparing two separate confidence intervals when the
  relevant object is a difference?
  \item Did I interpret a p-value as $\Prob(\text{data as extreme}\mid H_0)$,
  not as $\Prob(H_0\mid\text{data})$?
  \item Did I separate statistical significance from economic importance?
  \item Did I avoid saying ``no effect'' when the result is merely imprecise?
  \item Did I check whether many outcomes, subgroups, or specifications were
  searched?
  \item Did I explain the identifying assumption separately from the standard
  error?
\end{enumerate}

\section{Practice Questions}

\begin{enumerate}
  \item A sample has $n=25$, $\bar{Y}=10$, and $s=5$. Compute the standard
  error and a 95 percent confidence interval for the mean.
  \item If the sample size rises from $50$ to $200$ and the individual standard
  deviation is unchanged, by what factor does the standard error change?
  \item Explain why small schools may appear more often among both the best and
  worst schools in a ranking based on average test scores.
  \item In a randomized experiment, $\bar{Y}_1=84$, $\bar{Y}_0=80$,
  $s_1=s_0=12$, and $n_1=n_0=36$. Compute the treatment effect estimate,
  standard error, 95 percent confidence interval, and $z$ statistic for testing
  zero.
  \item Why is ``the p-value is 0.03, so there is a 3 percent probability that
  the null is true'' incorrect?
  \item Give an example of an estimate that is statistically significant but
  economically unimportant.
  \item Give an example of an estimate that is economically important but not
  statistically significant.
  \item What is power? Why do small studies often have low power?
  \item A paper reports 30 subgroup regressions and highlights the two with
  $p<0.05$. What concern should you raise?
  \item Explain in two sentences why randomization solves a selection problem
  but does not eliminate sampling variation.
\end{enumerate}

\end{document}
